% main.tex
\documentclass{article}

% VerifAI asks for ICLR format
\usepackage{iclr2026_conference,times}

\usepackage{microtype}
\usepackage{graphicx}
\usepackage{booktabs}
\usepackage{amsmath,amssymb}
\usepackage{url}
\usepackage{hyperref}

\title{NLVerifier: Natural Language Conditioned Cross-Encoder Selection for Lean Autoformalization Candidates}

\author{Anonymous authors\\
Paper under double-blind review
}

\begin{document}
\maketitle

\begin{abstract}
Sampling multiple candidate formalizations and selecting a final prediction is a common strategy in statement autoformalization.
Prior work uses type-check filtering followed by a surface-form consensus heuristic such as Self-BLEU.
We study a simple alternative: a learned natural language conditioned verifier (NLVerifier) that scores each (natural language statement, candidate Lean statement) pair and selects the highest scoring candidate.
On ProofNetVerif, NLVerifier improves top-1 selection accuracy from 84/178 (47.2\%) to 111/178 (62.4\%) on the original (transductive) test split, and from 17/55 (30.9\%) to 27/55 (49.1\%) on a new ID-disjoint (inductive) split constructed to prevent problem overlap.
Ablations that remove or collapse the natural language input reduce performance, indicating that the gains come from conditioning on the informal statement rather than only modeling Lean syntax.
\end{abstract}

\section{Introduction}
Statement autoformalization translates informal mathematical statements into formal languages such as Lean, enabling machine checking and downstream verification.
Recent evaluations emphasize that correctness cannot be reliably measured by surface overlap metrics alone, motivating equivalence-centric benchmarks and metrics \citep{poiroux2025reliable}.
In practical pipelines, models often sample many candidate formalizations, discard those that fail to type-check, then select a single prediction using a heuristic such as Self-BLEU \citep{poiroux2025reliable,zhu2018texygen}.
These heuristics are attractive because they require no training, but they ignore the natural language prompt and can fail under distribution shift.

This paper evaluates a small, direct modification: replace the consensus heuristic with a learned verifier that explicitly scores alignment between the natural language statement and each candidate Lean statement.
Our goal is not to propose a new full autoformalization system, but to isolate and strengthen the selection component that sits between sampling and final evaluation.

\section{Task}
We consider the following selection setting.
For each problem ID $p$, we observe a natural language statement $x_p$ and a set of $n_p$ candidate Lean statements $\{y_{p,i}\}_{i=1}^{n_p}$.
Each candidate is annotated with a binary semantic equivalence label $c_{p,i} \in \{0,1\}$ indicating whether it is correct for $x_p$.
A selection method outputs a single index $\hat{i}$ and is scored by whether $c_{p,\hat{i}} = 1$.

\paragraph{Top-1 selection accuracy}
We report \texttt{selected\_correct}, the fraction of problems where the selected candidate is correct.

\paragraph{Conditional accuracy}
We also report \texttt{selected\_correct\_given\_any}, the fraction of problems where the method selects a correct candidate conditional on the candidate set containing at least one correct option.
This isolates selection quality from candidate coverage.

\section{Methods}

\subsection{Self-BLEU consensus baseline}
Self-BLEU was introduced as a diversity diagnostic \citep{zhu2018texygen}, and has been used as a practical selection heuristic in autoformalization pipelines \citep{poiroux2025reliable}.
We use it as a surface-form consensus selector by choosing the candidate with the highest average BLEU score against the other candidates for the same problem.

\subsection{Random baseline}
As a lower bound, we select a candidate uniformly at random from each candidate set.

\subsection{Candidate-candidate critic baseline}
We include a learned candidate-only baseline that scores candidates without access to the natural language input.
This can be implemented as a cross-encoder that predicts semantic compatibility between two Lean candidates and selects a medoid under the learned similarity.
It captures stronger Lean-aware similarity than BLEU, but it cannot directly test alignment with the informal statement.

\subsection{NLVerifier}
NLVerifier is a cross-encoder reranker \citep{nogueira2019passage} that produces a scalar score $s_\theta(x_p, y_{p,i})$ for each (natural language, Lean) pair.
We instantiate NLVerifier with DeBERTa-v3-base \citep{he2021debertav3} and fine-tune it as a binary classifier on labeled pairs.
At inference time, NLVerifier selects:
\[
\hat{i} = \arg\max_i s_\theta(x_p, y_{p,i}).
\]

\paragraph{Ablations}
To test whether the model truly uses the natural language input, we evaluate two inference-time ablations:
(1) \textbf{blank NL}, where $x_p$ is replaced by an empty string, and
(2) \textbf{constant NL}, where $x_p$ is replaced by the same constant string for all problems.

\section{Evaluation setup}

\subsection{ProofNetVerif}
ProofNetVerif is a benchmark introduced to evaluate autoformalization metrics and equivalence labeling \citep{poiroux2025reliable,proofnetverif_hf}.
We use the dataset in a local packaging that provides an ``original'' split and a derived ID-disjoint split.

\paragraph{Candidate sets and construction}
ProofNetVerif is stored as a flat table of candidate rows keyed by a problem ID.
Each row contains a natural language statement, a candidate Lean statement, and a binary equivalence label.
For selection evaluation, we construct one candidate set per problem by grouping rows with the same problem ID and taking each row's candidate Lean statement as one selectable option.

On the transductive test split, we evaluate 178 problems with 1{,}452 total candidates (mean 8.16 candidates per problem), and 120/178 problems (67.4\%) contain at least one correct candidate.
On the ID-disjoint inductive test split, we evaluate 55 problems with 542 total candidates (mean 9.85 candidates per problem), and 29/55 problems (52.7\%) contain at least one correct candidate.

\begin{table}[t]
\centering
\small
\caption{Candidate set statistics for evaluation splits.}
\label{tab:candstats}
\begin{tabular}{lrrrr}
\toprule
Split & Problems & Candidates & Mean cand./prob. & Any-correct problems \\
\midrule
Transductive test & 178 & 1{,}452 & 8.16 & 120 (67.4\%) \\
ID-disjoint test & 55 & 542 & 9.85 & 29 (52.7\%) \\
\bottomrule
\end{tabular}
\end{table}

\paragraph{Transductive evaluation}
We evaluate on the original test split, which is transductive with respect to the dataset's original split design.

\paragraph{Inductive evaluation via ID-disjoint splitting}
To test generalization to unseen problems, we construct a new split by grouping examples by problem ID and ensuring IDs do not overlap between train, dev, and test.
Our ID-disjoint split yields 252 train IDs (2{,}724 rows), 54 dev IDs (486 rows), and 55 test IDs (542 rows).
We report inductive results on the 55-ID test partition.

\section{Results}

\IfFileExists{paper/generated/nlverifier_main_table.tex}{
  % Generated by scripts/summarize_nlverifier_paper_metrics.py; do not edit by hand.
\begin{table}[t]
\centering
\small
\caption{Selection accuracy on ProofNetVerif, reported as correct/total (percent).
``Given-any'' conditions on problems whose candidate set contains at least one correct candidate
(120/178 transductive, 29/55 inductive).}
\label{tab:main}
\resizebox{\textwidth}{!}{%
\begin{tabular}{lcccc}
\toprule
& \multicolumn{2}{c}{Transductive (178 problems)} & \multicolumn{2}{c}{Inductive ID-disjoint (55 problems)} \\
Method & Top-1 & Given-any & Top-1 & Given-any \\
\midrule
Random selection &
63/178 (35.4\%) & 63/120 (52.5\%) & 12/55 (21.8\%) & 12/29 (41.4\%) \\
Self-BLEU consensus \citep{zhu2018texygen,poiroux2025reliable} &
84/178 (47.2\%) & 84/120 (70.0\%) & 17/55 (30.9\%) & 17/29 (58.6\%) \\
Candidate-only critic (best variant) &
87/178 (48.9\%) & 87/120 (72.5\%) & 20/55 (36.4\%) & 20/29 (69.0\%) \\
\midrule
NLVerifier (full input) &
\textbf{111/178 (62.4\%)} & \textbf{111/120 (92.5\%)} & \textbf{27/55 (49.1\%)} & \textbf{27/29 (93.1\%)} \\
NLVerifier (blank NL) &
101/178 (56.7\%) & 101/120 (84.2\%) & 23/55 (41.8\%) & 23/29 (79.3\%) \\
NLVerifier (constant NL) &
104/178 (58.4\%) & 104/120 (86.7\%) & 25/55 (45.5\%) & 25/29 (86.2\%) \\
\bottomrule
\end{tabular}
}
\end{table}

}{
  % Generated by scripts/summarize_nlverifier_paper_metrics.py; do not edit by hand.
\begin{table}[t]
\centering
\small
\caption{Selection accuracy on ProofNetVerif, reported as correct/total (percent).
``Given-any'' conditions on problems whose candidate set contains at least one correct candidate
(120/178 transductive, 29/55 inductive).}
\label{tab:main}
\resizebox{\textwidth}{!}{%
\begin{tabular}{lcccc}
\toprule
& \multicolumn{2}{c}{Transductive (178 problems)} & \multicolumn{2}{c}{Inductive ID-disjoint (55 problems)} \\
Method & Top-1 & Given-any & Top-1 & Given-any \\
\midrule
Random selection &
63/178 (35.4\%) & 63/120 (52.5\%) & 12/55 (21.8\%) & 12/29 (41.4\%) \\
Self-BLEU consensus \citep{zhu2018texygen,poiroux2025reliable} &
84/178 (47.2\%) & 84/120 (70.0\%) & 17/55 (30.9\%) & 17/29 (58.6\%) \\
Candidate-only critic (best variant) &
87/178 (48.9\%) & 87/120 (72.5\%) & 20/55 (36.4\%) & 20/29 (69.0\%) \\
\midrule
NLVerifier (full input) &
\textbf{111/178 (62.4\%)} & \textbf{111/120 (92.5\%)} & \textbf{27/55 (49.1\%)} & \textbf{27/29 (93.1\%)} \\
NLVerifier (blank NL) &
101/178 (56.7\%) & 101/120 (84.2\%) & 23/55 (41.8\%) & 23/29 (79.3\%) \\
NLVerifier (constant NL) &
104/178 (58.4\%) & 104/120 (86.7\%) & 25/55 (45.5\%) & 25/29 (86.2\%) \\
\bottomrule
\end{tabular}
}
\end{table}

}

Table~\ref{tab:main} shows that NLVerifier improves top-1 selection accuracy over Self-BLEU
from 84/178 (47.2\%) to 111/178 (62.4\%) on the transductive test split, and from 17/55 (30.9\%) to 27/55 (49.1\%) on the ID-disjoint inductive split.
When at least one correct candidate exists in the candidate set, NLVerifier selects a correct candidate
in 111/120 cases (92.5\%) transductively and 27/29 cases (93.1\%) inductively.

\paragraph{Does NL conditioning matter}
Both ablations reduce performance, especially on the inductive split, supporting that NLVerifier uses the natural language statement and is not only learning Lean plausibility priors.

\section{Limitations and future work}
This study focuses on one primary selection benchmark.
As a preliminary stress test, we also evaluated zero-shot transfer to an external misalignment dataset and observed limited transfer, which we summarize in Appendix~\ref{app:ood}.
Two natural next steps are (1) training on a mixture of alignment and equivalence datasets spanning multiple benchmarks, and (2) integrating the verifier into an end-to-end autoformalization pipeline to measure downstream gains in statement correctness after type-check filtering.

\section{Conclusion}
We presented NLVerifier, a simple natural language conditioned cross-encoder for selecting among sampled Lean autoformalization candidates.
On ProofNetVerif, NLVerifier yields large gains over Self-BLEU and remains substantially stronger under an ID-disjoint inductive evaluation.

\section*{Reproducibility}
For double-blind review, we provide code, splits, and all paper tables in supplementary material via an anonymized repository and artifact bundle.

\bibliography{references}
\bibliographystyle{iclr2026_conference}

\appendix

\section{Out-of-distribution stress test}
\label{app:ood}
A concern with any learned verifier is generalization beyond the training distribution.
As a stress test, we applied NLVerifier zero-shot to a misalignment benchmark from FormalAlign \citep{lu2024formalalign} (MiniF2F misalignment test).
FormalAlign constructs adversarially misaligned formal statements and evaluates both detection and selection settings.

In our zero-shot evaluation, we observed 73.6\% overall pair classification accuracy at a fixed threshold, with lower accuracy on the positive class (26.3\%) than on the negative class (76.7\%).
When cast as per-input selection among candidates, top-1 selection accuracy was 2.47\%.
These results indicate substantial domain shift between ProofNetVerif equivalence labels and FormalAlign's adversarial misalignment construction, motivating multi-benchmark training as future work.

\end{document}
